\chapter{Solutions to Chapter 26: Principal Stratification}\label{sol:chapter26}

\begin{exercisesolution}{Complete the proof of \cref{thm::identification-PCEs-pscore}}{strong monotonicity 下 $\tau(0,0)$ 的 principal score 识别}
本题补全 \cref{thm::identification-PCEs-pscore} 中第二个 principal causal effect 的证明。与正文已经证明的 $\tau(1,0)$ 一样，关键不是直接观察 principal stratum，而是用 principal score 在 control group 中重构该 stratum 的 $Y(0)$ 均值。

在 \cref{assume::strong-monotonicity} 下，$M(0)=0$，所以 principal stratum 只由 $M(1)$ 决定：
\[
M(1)=1 \quad \text{对应 stratum } (1,0),
\qquad
M(1)=0 \quad \text{对应 stratum } (0,0).
\]
由 \cref{assume::complete-randomization-withM} 和 consistency，
\[
\pi(X)=\mathbb{P}\{M(1)=1\mid X\}
=\mathbb{P}(M=1\mid Z=1,X),
\]
且
\[
\pi=\mathbb{P}\{M(1)=1\}=\mathbb{P}(M=1\mid Z=1).
\]
因此 $1-\pi(X)=\mathbb{P}\{M(1)=0\mid X\}$，$1-\pi=\mathbb{P}\{M(1)=0\}$。

先看 treated potential outcome 的部分。由于 $Z=1,M=0$ 等价于 $Z=1,M(1)=0$，随机化给出
\[
E(Y\mid Z=1,M=0)
=E\{Y(1)\mid Z=1,M(1)=0\}
=E\{Y(1)\mid M(1)=0\}.
\]
这正是 $\tau(0,0)$ 中的 treated mean。

接下来识别 control potential outcome 的部分。我们要证明
\[
\frac{E\{(1-\pi(X))Y\mid Z=0\}}{1-\pi}
=E\{Y(0)\mid M(1)=0\}.
\]
由 consistency 与 \cref{assume::complete-randomization-withM}，
\[
E\{(1-\pi(X))Y\mid Z=0\}
=E\{(1-\pi(X))Y(0)\}.
\]
再对 $X$ 条件化，并使用 \cref{assume::principal-ignorability-SM} 所蕴含的
\[
E\{Y(0)\mid M(1),X\}=E\{Y(0)\mid X\},
\]
得到
\begin{align*}
E\{(1-\pi(X))Y(0)\}
&=
E\left[(1-\pi(X))E\{Y(0)\mid X\}\right]\\
&=
E\left[\mathbb{P}\{M(1)=0\mid X\}E\{Y(0)\mid X\}\right]\\
&=
E\left[E\{(1-M(1))Y(0)\mid X\}\right]\\
&=
E\{(1-M(1))Y(0)\}.
\end{align*}
两边除以 $1-\pi=\mathbb{P}\{M(1)=0\}$，即得
\[
\frac{E\{(1-\pi(X))Y\mid Z=0\}}{1-\pi}
=E\{Y(0)\mid M(1)=0\}.
\]
于是
\[
\tau(0,0)
=E\{Y(1)-Y(0)\mid M(1)=0\}
=E(Y\mid Z=1,M=0)
-\frac{E\{(1-\pi(X))Y\mid Z=0\}}{1-\pi},
\]
这证明了 \cref{thm::identification-PCEs-pscore} 中关于 $\tau(0,0)$ 的公式。

\paragraph{Remark.}
这道题说明 principal score 的作用与 propensity score 平行：它把一个 latent principal stratum 的分布信息压缩为一个 conditional probability。不同之处在于，propensity score 平衡 treatment assignment，而 principal score 平衡不可直接观测的 principal strata。
\end{exercisesolution}

\begin{exercisesolution}{Principal score method under monotonicity}{一般 monotonicity 下三个 principal causal effects 的识别}
本题把 \cref{thm::identification-PCEs-pscore} 从 strong monotonicity 推广到 ordinary monotonicity。此时 $M(1)\geq M(0)$，所以 latent principal strata 为
\[
(1,1),\qquad (1,0),\qquad (0,0),
\]
分别对应 always survivors、protected units 和 never survivors。困难在于 $(Z=1,M=1)$ 是 $(1,1)$ 与 $(1,0)$ 的 mixture，而 $(Z=0,M=0)$ 是 $(1,0)$ 与 $(0,0)$ 的 mixture。权重 $w_{z,u}(X)$ 的作用正是从这些 observed mixtures 中抽取目标 principal stratum。

先证明 principal scores 的识别。由 \cref{assume::complete-randomization-withM}、consistency 与 \cref{assume::monotonicity-ps}，
\[
\mathbb{P}(M=0\mid Z=1,X)
=\mathbb{P}\{M(1)=0\mid X\}
=\pi_{(0,0)}(X),
\]
\[
\mathbb{P}(M=1\mid Z=0,X)
=\mathbb{P}\{M(0)=1\mid X\}
=\pi_{(1,1)}(X),
\]
并且
\begin{align*}
\mathbb{P}(M=1\mid Z=1,X)-\mathbb{P}(M=1\mid Z=0,X)
&=\mathbb{P}\{M(1)=1\mid X\}-\mathbb{P}\{M(0)=1\mid X\}\\
&=\pi_{(1,0)}(X).
\end{align*}
去掉对 $X$ 的条件即可得到 marginal principal scores 的三个公式。

下面证明 causal effects 的识别。记 $U=(M(1),M(0))$。我们先处理 $\tau(1,0)$ 的 treated mean。令
\[
r_{1,(1,0)}(X)
=\frac{\pi_{(1,0)}(X)}{\pi_{(1,0)}(X)+\pi_{(1,1)}(X)}.
\]
它是在给定 $X$ 且 $M(1)=1$ 时，unit 属于 stratum $(1,0)$ 的 conditional probability。由 \cref{assume::principal-ignorability-M} 的第一部分，
\[
E\{Y(1)\mid U=(1,0),X\}
=E\{Y(1)\mid U=(1,1),X\}.
\]
因此给定 $X$，在 observed mixture $Z=1,M=1$ 中，$Y(1)$ 的 conditional mean 不随 $(1,0)$ 与 $(1,1)$ 的 stratum membership 改变。于是
\begin{align*}
&E\{w_{1,(1,0)}(X)Y\mid Z=1,M=1\}\\
&\quad =
\frac{\pi_{(1,0)}+\pi_{(1,1)}}{\pi_{(1,0)}}
E\{r_{1,(1,0)}(X)Y(1)\mid M(1)=1\}\\
&\quad =
\frac{1}{\pi_{(1,0)}}
E\left[\pi_{(1,0)}(X)E\{Y(1)\mid U=(1,0),X\}\right]\\
&\quad =
E\{Y(1)\mid U=(1,0)\}.
\end{align*}
同理，令
\[
r_{0,(1,0)}(X)
=\frac{\pi_{(1,0)}(X)}{\pi_{(1,0)}(X)+\pi_{(0,0)}(X)}.
\]
由 \cref{assume::principal-ignorability-M} 的第二部分，
\[
E\{Y(0)\mid U=(1,0),X\}
=E\{Y(0)\mid U=(0,0),X\}.
\]
在 observed mixture $Z=0,M=0$ 中加权得到
\[
E\{w_{0,(1,0)}(X)Y\mid Z=0,M=0\}
=E\{Y(0)\mid U=(1,0)\}.
\]
两式相减，即得 \cref{thm::identification-PCEs-pscore-general-ding-lu} 中 $\tau(1,0)$ 的公式。

再看 $\tau(0,0)$。当 $Z=1,M=0$ 时，由 monotonicity 必有 $U=(0,0)$，所以
\[
E(Y\mid Z=1,M=0)=E\{Y(1)\mid U=(0,0)\}.
\]
另一方面，在 $Z=0,M=0$ 的 mixture 中使用
\[
r_{0,(0,0)}(X)
=\frac{\pi_{(0,0)}(X)}{\pi_{(1,0)}(X)+\pi_{(0,0)}(X)}
\]
并重复上面的条件化论证，得到
\[
E\{w_{0,(0,0)}(X)Y\mid Z=0,M=0\}
=E\{Y(0)\mid U=(0,0)\}.
\]
这给出 $\tau(0,0)$ 的识别公式。

最后看 $\tau(1,1)$。当 $Z=0,M=1$ 时，由 monotonicity 必有 $U=(1,1)$，所以
\[
E(Y\mid Z=0,M=1)=E\{Y(0)\mid U=(1,1)\}.
\]
在 $Z=1,M=1$ 的 mixture 中使用
\[
r_{1,(1,1)}(X)
=\frac{\pi_{(1,1)}(X)}{\pi_{(1,0)}(X)+\pi_{(1,1)}(X)}
\]
并利用 \cref{assume::principal-ignorability-M} 的第一部分，得到
\[
E\{w_{1,(1,1)}(X)Y\mid Z=1,M=1\}
=E\{Y(1)\mid U=(1,1)\}.
\]
两式相减，即得 $\tau(1,1)$ 的识别公式。

\paragraph{Remark.}
这一定理的形式看起来比 \cref{thm::identification-PCEs-pscore} 复杂，但底层逻辑没有变：pure observed cells 直接给出某些 stratum-specific means，mixture cells 则通过 principal score weights 拆开。principal ignorability 是让这种拆分可以只依赖 $X$ 的核心不可检验假设。
\end{exercisesolution}

\begin{exercisesolution}{Principal score method in observational studies}{observational studies 中 strong monotonicity 版本的 IPW 识别}
本题把 \cref{thm::identification-PCEs-pscore} 从 randomized experiment 推广到 observational study。新的难点是 $Z$ 不再边际随机化，而只在给定 $X$ 后满足 \cref{assume::ignorabilitywithM}。因此除了 principal score $\pi(X)$，还必须使用 propensity score
\[
e(X)=\mathbb{P}(Z=1\mid X)
\]
来校正 treatment assignment 的选择偏差。下面默认 standard positivity：$0<e(X)<1$，且 $0<\pi<1$。

先证明 principal score 的识别。在 \cref{assume::strong-monotonicity} 下，$M=M(1)$ when $Z=1$。由 consistency 与 \cref{assume::ignorabilitywithM}，
\[
\mathbb{P}(M=1\mid Z=1,X)
=\mathbb{P}\{M(1)=1\mid Z=1,X\}
=\mathbb{P}\{M(1)=1\mid X\}
=\pi(X).
\]
因此
\[
\pi=E\{\pi(X)\}=E\{\mathbb{P}(M=1\mid Z=1,X)\}.
\]
注意这里一般不能把 $\pi$ 写成 $\mathbb{P}(M=1\mid Z=1)$，因为 observational treated group 的 $X$ 分布不等于总体 $X$ 分布。

接下来证明 $\tau(1,0)$ 的公式。由 inverse probability weighting，
\begin{align*}
E\left\{\frac{M}{\pi}\frac{Z}{e(X)}Y\right\}
&=
\frac{1}{\pi}
E\left[E\left\{\frac{Z}{e(X)}M Y\mid X,M(1),Y(1)\right\}\right]\\
&=
\frac{1}{\pi}E\{M(1)Y(1)\}\\
&=
E\{Y(1)\mid M(1)=1\}.
\end{align*}
这里第一步使用 observed data 中 $Z=1$ 时 $M=M(1)$ 且 $Y=Y(1)$。对于 control potential outcome，
\begin{align*}
E\left\{\frac{\pi(X)}{\pi}\frac{1-Z}{1-e(X)}Y\right\}
&=
\frac{1}{\pi}E\{\pi(X)Y(0)\}\\
&=
\frac{1}{\pi}E\left[\pi(X)E\{Y(0)\mid X\}\right]\\
&=
\frac{1}{\pi}E\left[E\{M(1)Y(0)\mid X\}\right]\\
&=
E\{Y(0)\mid M(1)=1\},
\end{align*}
其中第三步使用 \cref{assume::principal-ignorability-SM}。两式相减得到 $\tau(1,0)$ 的识别公式。

对 $\tau(0,0)$ 完全类似。treated mean 为
\begin{align*}
E\left\{\frac{1-M}{1-\pi}\frac{Z}{e(X)}Y\right\}
&=
\frac{1}{1-\pi}E\{(1-M(1))Y(1)\}\\
&=
E\{Y(1)\mid M(1)=0\}.
\end{align*}
control mean 为
\begin{align*}
E\left\{\frac{1-\pi(X)}{1-\pi}\frac{1-Z}{1-e(X)}Y\right\}
&=
\frac{1}{1-\pi}E\{(1-\pi(X))Y(0)\}\\
&=
\frac{1}{1-\pi}E\{(1-M(1))Y(0)\}\\
&=
E\{Y(0)\mid M(1)=0\},
\end{align*}
其中第二步仍然由 \cref{assume::principal-ignorability-SM} 推出。于是
\[
\tau(0,0)
=E\left\{\frac{1-M}{1-\pi}\frac{Z}{e(X)}Y\right\}
-E\left\{\frac{1-\pi(X)}{1-\pi}\frac{1-Z}{1-e(X)}Y\right\}.
\]
这完成了 \cref{thm::identification-PCEs-pscore-obs} 的证明。

\paragraph{Remark.}
observational study 版本有两个权重层次：$Z/e(X)$ 与 $(1-Z)/(1-e(X))$ 负责把 observational data 变回目标总体下的 potential outcomes；$\pi(X)$ 与 $1-\pi(X)$ 负责把总体中的 control potential outcomes 分配到 latent principal strata。
\end{exercisesolution}

\begin{exercisesolution}{General principal score method}{observational studies 中 monotonicity 版本的统一公式}
本题要求把 \cref{thm::identification-PCEs-pscore-general-ding-lu} 与 \cref{thm::identification-PCEs-pscore-obs} 合并。也就是说，我们同时允许 ordinary monotonicity 下的三个 principal strata，并允许 treatment assignment 只在给定 $X$ 后满足 ignorability。最终公式就是：在 randomized theorem 的 principal-score weights 外，再乘上 propensity-score 的 inverse probability weights。

令
\[
U=(M(1),M(0))\in\{(1,1),(1,0),(0,0)\},
\]
并定义
\[
\pi_u(X)=\mathbb{P}(U=u\mid X),
\qquad
\pi_u=E\{\pi_u(X)\}.
\]
在 \cref{assume::ignorabilitywithM,assume::monotonicity-ps} 下，
\begin{align*}
\pi_{(0,0)}(X)&=\mathbb{P}(M=0\mid Z=1,X),\\
\pi_{(1,1)}(X)&=\mathbb{P}(M=1\mid Z=0,X),\\
\pi_{(1,0)}(X)&=\mathbb{P}(M=1\mid Z=1,X)-\mathbb{P}(M=1\mid Z=0,X),
\end{align*}
且 marginal principal scores 由 $\pi_u=E\{\pi_u(X)\}$ 识别。

在 standard positivity 条件下，三个 principal causal effects 可由下面的 formulas 识别：
\begin{align*}
\tau(1,0)
&=
\frac{1}{\pi_{(1,0)}}E\left\{
\frac{Z}{e(X)}M
\frac{\pi_{(1,0)}(X)}{\pi_{(1,0)}(X)+\pi_{(1,1)}(X)}
Y\right\}\\
&\quad -
\frac{1}{\pi_{(1,0)}}E\left\{
\frac{1-Z}{1-e(X)}(1-M)
\frac{\pi_{(1,0)}(X)}{\pi_{(1,0)}(X)+\pi_{(0,0)}(X)}
Y\right\},
\end{align*}
\begin{align*}
\tau(0,0)
&=
\frac{1}{\pi_{(0,0)}}E\left\{
\frac{Z}{e(X)}(1-M)Y\right\}\\
&\quad -
\frac{1}{\pi_{(0,0)}}E\left\{
\frac{1-Z}{1-e(X)}(1-M)
\frac{\pi_{(0,0)}(X)}{\pi_{(1,0)}(X)+\pi_{(0,0)}(X)}
Y\right\},
\end{align*}
以及
\begin{align*}
\tau(1,1)
&=
\frac{1}{\pi_{(1,1)}}E\left\{
\frac{Z}{e(X)}M
\frac{\pi_{(1,1)}(X)}{\pi_{(1,0)}(X)+\pi_{(1,1)}(X)}
Y\right\}\\
&\quad -
\frac{1}{\pi_{(1,1)}}E\left\{
\frac{1-Z}{1-e(X)}MY\right\}.
\end{align*}

下面证明这些公式。先看 $\tau(1,0)$ 的 treated mean。由于 $Z=1,M=1$ 对应 $U=(1,0)$ 或 $U=(1,1)$，而 \cref{assume::principal-ignorability-M} 的第一部分使 $Y(1)$ 在这两个 strata 中给定 $X$ 后具有相同 conditional mean，
\begin{align*}
&E\left\{
\frac{Z}{e(X)}M
\frac{\pi_{(1,0)}(X)}{\pi_{(1,0)}(X)+\pi_{(1,1)}(X)}
Y\right\}\\
&\quad =
E\left[\pi_{(1,0)}(X)E\{Y(1)\mid U=(1,0),X\}\right]\\
&\quad =
E\{1(U=(1,0))Y(1)\}.
\end{align*}
除以 $\pi_{(1,0)}$ 得到 $E\{Y(1)\mid U=(1,0)\}$。同理，使用 \cref{assume::principal-ignorability-M} 的第二部分，
\begin{align*}
&E\left\{
\frac{1-Z}{1-e(X)}(1-M)
\frac{\pi_{(1,0)}(X)}{\pi_{(1,0)}(X)+\pi_{(0,0)}(X)}
Y\right\}\\
&\quad =
E\{1(U=(1,0))Y(0)\}.
\end{align*}
除以 $\pi_{(1,0)}$ 后得到 $E\{Y(0)\mid U=(1,0)\}$，故 $\tau(1,0)$ 公式成立。

对 $\tau(0,0)$，treated side 是 pure cell：
\[
E\left\{\frac{Z}{e(X)}(1-M)Y\right\}
=E\{1(U=(0,0))Y(1)\}.
\]
control side 位于 mixture cell $Z=0,M=0$ 中；由 \cref{assume::principal-ignorability-M} 的第二部分，
\[
E\left\{
\frac{1-Z}{1-e(X)}(1-M)
\frac{\pi_{(0,0)}(X)}{\pi_{(1,0)}(X)+\pi_{(0,0)}(X)}
Y\right\}
=E\{1(U=(0,0))Y(0)\}.
\]
分别除以 $\pi_{(0,0)}$ 并相减，得到 $\tau(0,0)$。

对 $\tau(1,1)$，control side 是 pure cell：
\[
E\left\{\frac{1-Z}{1-e(X)}MY\right\}
=E\{1(U=(1,1))Y(0)\}.
\]
treated side 位于 mixture cell $Z=1,M=1$ 中；由 \cref{assume::principal-ignorability-M} 的第一部分，
\[
E\left\{
\frac{Z}{e(X)}M
\frac{\pi_{(1,1)}(X)}{\pi_{(1,0)}(X)+\pi_{(1,1)}(X)}
Y\right\}
=E\{1(U=(1,1))Y(1)\}.
\]
分别除以 $\pi_{(1,1)}$ 并相减，即得 $\tau(1,1)$。

\paragraph{Remark.}
这道题给出的是 principal score 方法的统一模板。随机实验中，$e(X)$ 是已知常数，公式会退化为 \cref{thm::identification-PCEs-pscore-general-ding-lu} 的条件均值版本；strong monotonicity 下，$(1,1)$ stratum 消失，公式又退化为 \cref{thm::identification-PCEs-pscore-obs}。
\end{exercisesolution}
